\section{Freedom and Destiny}

\paragraph{Multiple Selves}


We saw from Valdimir Solovyov's description that, in the second moment or positing, the divine Unity is objectified as multiplicity. Specifically, multiple selves arise out of the void. So what was in the divine mind as undifferentiated possibilities manifests as differentiated actualities.

These are free beings with all that implies. Each such being has its possibilities and its destiny is to manifest those possibilities. This is the metaphysical, esoteric, and Hermetic understanding of creation and is consistent, even if expressed slightly differently, with Rene Guenon, Vladimir Solovyov, and Valentin Tomberg. It accounts for our experience of a ``previous'' existence, with priority understood ontologically, not temporally. Keep in mind, too, that each such being comes into existence in the physical and psychical environment suitable to it. This can explain our feeling, or remembrance, of ``being there'', especially the psychic current. From that we get our feeling of at-homeness: identity, continuity, community.

Unfortunately, in our time, the so-called Kali Yuga, that sense of dwelling is often lacking and many are lost in a feeling of alienation. The four elements identified by Martin Heidegger as creating dwelling---earth, sky, people, the gods---are opaque to some extent. They are not a ``given'', but must be (re)created, the destiny of those who understand this.

\paragraph{Freedom and Multiplicity}


In the Emperor Arcanum, Tomberg develops the idea of freedom. In an analogous way to God's three positings, each man recapitulates the same three moments. Man, too, is free. His existence is freedom, i.e., the void within him, analogous to the idea of the Void. His essence, then, is the fullness of his possibilities. Unlike God, whose existence and essence coincide, man's task is to manifest his possibilities in freedom.

This requires a conscious participation in the three moments, in an analogous way:

\begin{enumerate}


\item He must recall his destiny, i.e., what his possibilities are, in the specific world he finds himself. In the fall into multiplicity, there is a forgetfulness, since each created being embodies an empirical self, i.e., one conditioned in space and time. Shankara, for example, calls these conditionings Upadhis. The being will take this self as his real self, thereby forgetting his original purpose. For man, then, freedom is not an absolute, but is conditioned by external and internal factors. To be completely free, he must learn to detach himself from all such conditionings.

\item Actually, the being does not project or objectify a single empirical self, but man has many such selves, analogous to God's second positing. This is seldom noticed in our daily consciousness, as each empirical self tries to claim authority. It is easy to recognize that we create all the characters in the dream state, while identifying with just one of them. However, in the waking state, we fail to see the same multitude and falsely believe we are one. Hence, a spiritual practice must involve overcoming such an illusion. One's multiplicity needs to be experienced directly, otherwise there will be no motivation to proceed to the third positing.

\item In the third positing, true unity is restored. That is, there develops and awareness of one's real self, the noumenal self above space and time. Only in that state can one manifest his destiny.
\end{enumerate}


\paragraph{Salvation and Freedom}


Some may interpret the idea that different men are born with various gifts and possibilities as an injustice. That cannot be. Possibilities are just what they are, and they distinguish people from each other, making each one unique. The alternative is a type of egalitarian regimentation, such as the images we get of crowds in North Korea. That is the extreme case, but regimentation of thought and opinion can be found everywhere.

The doctrine of karma and reincarnation are supposed to account for that alleged injustice, but they fail in that for reasons already given. Rather than elevating the moral life, it provides the lowest level of motivation, to wit, that of naked self-interest.

Furthermore, the idea of karma is actually misguided. A man may be born into a life of leisure for ``good'' deeds in a past life. But is such a lazy life causes him to forget and neglect his true destiny, then how can that be accounted as a good?

No, what this really means is that salvation or liberation is our real destiny, and not merely the amelioration of the material circumstances of life. In this sense, salvation is potentially available to all, regardless of circumstances of birth. Hence there is no injustice in creation.

In this understanding, we are free to choose our post-mortem destiny; this is certainly the most just teaching. There is the oft-heard complaint that it is unjust because there is uncertainty about it. They want signs and wonders. Paradoxically, they want to be ``compelled'' to act, when freedom is the opposite of compulsion. A serious man takes his life seriously. Once he understands his true situation, then he will compel himself to work out his own salvation.


\flrightit{Posted on 2014-05-18 by Cologero}

